% ABSTRACT (draft v6). Pure Rust causal anomaly engine on DeepCausality.
% MAIN RESULT = the SYSTEM: a real-time, edge-resident anomaly engine across a
% fleet of system/network monitoring agents, enabled by the DeepCausality reasoner.
% Accuracy fix: the shipped edge path is the STATELESS one (keep window, recompute
% Welford each eval), not an O(1) running accumulator. Numbers from a fresh
% Apple M4 run (~0.5M evals/s/core warm).
% Humanizer: no em/en dashes; plain technical voice; simple copulas.
\sysname runs system and network monitoring across a fleet of agents, and those
agents emit millions of metric time series. Catching anomalous metrics across the
fleet, in real time, is the problem this paper addresses. We present an
anomaly-detection engine that runs at the edge, on each agent, built on \dc{}, a
high-performance computational-causality framework. Each metric series is watched
by a \emph{causaloid}, and the detector's control flow is written as a \dc{}
\texttt{CausalFlow} pipeline. Framing each series as a causaloid keeps the detector
small and composable, lets the same core produce identical verdicts on the agent
and in the core, and leaves open a path to causal interventions. The reasoner keeps
only a bounded window of recent
samples per series, computes the running mean and variance with Welford's method
over that window, and raises a verdict when a sample's $z$-score stays past a
threshold for several consecutive samples. Keeping just the window, rather than a
running accumulator, makes each series cheap to checkpoint and restart. Verdicts
are emitted as Open Cybersecurity Schema Framework findings. The engine is written
entirely in Rust and runs as a resource-capped add-on alongside each monitoring
agent. On a single core it sustains roughly half a million sample evaluations per
second, far more than an agent produces, so the engine fits inside a small CPU and
memory budget next to the agent and flags anomalies as the data is produced rather
than shipping every raw point to a central service. We describe the engine's
design, its detection algorithm, and microbenchmarks of the reasoner, and we are
explicit about what we have not yet measured, chiefly detection accuracy on
labeled data.
